\section{Performa: Praktik Baik dan Profiling}

Performa JavaScript mempengaruhi pengalaman pengguna, terutama di perangkat mobile. Bottleneck umum: DOM manipulation berlebihan (\textit{layout thrashing}), event handler tanpa debounce/throttle, objek besar yang disalin berulang, dan blocking long tasks yang membekukan UI \cite{mdnJsGuide}.

Praktik baik: (1) batch operasi DOM (gunakan \code{DocumentFragment}); (2) debounce/throttle event handler; (3) gunakan \code{requestAnimationFrame()} untuk animasi; (4) lazy load gambar dan modul; (5) profil kode dengan DevTools Performance tab untuk mengidentifikasi bottleneck; (6) gunakan \code{Web Workers} untuk komputasi berat agar tidak memblokir UI thread \cite{mdnToolsTesting}.

\begin{konsep}
\textbf{Layout Thrashing.} Membaca properti layout (\code{offsetHeight}, \code{getBoundingClientRect()}) lalu menulis ke DOM secara bergantian memaksa browser menghitung layout berulang kali (forced synchronous layout). Solusi: kumpulkan semua pembacaan, lalu lakukan semua penulisan. Library seperti \code{fastdom} membantu mengelola ini secara otomatis \cite{mdnJsGuide}.
\end{konsep}

